\begin{table}[b]
    \centering
    \begin{ruledtabular}
    \begin{tabular}{ccccc}
        $T$ & $L$ & $\tau_{2}$ & $\tau_{3}$ & $\tau_{4}$ \\\hline
         $0.7$ & $16$ & 1.445(5) & 1.730(11) & 1.941(16)\\
         $0.7$ & $24$ & 1.435(12) & 1.72(3) & 1.93(4)\\
         $0.8$ & $16$ & 1.448(5) & 1.732(11) & 1.943(17)\\
         $0.8$ & $24$ & 1.456(12) & 1.75(3) & 1.97(4)\\
         $0.9$ & $16$ & 1.450(5) & 1.736(11) & 1.941(19)\\
         $0.9$ & $24$ & 1.448(12) & 1.730(24) & 1.95(4)\\
         $1.0$ & $16$ & 1.463(6) & 1.759(13) & 1.968(22)\\
         $1.0$ & $24$ & 1.445(14) & 1.72(3) & 1.91(4)\\
         $T_\mathrm{c}$ & $16$ & 1.478(7) & 1.782(16) & 2.00(3)\\
         $T_\mathrm{c}$ & $24$ & 1.442(17) & 1.71(4) & 1.88(6)\\\hline
         \multicolumn{2}{c}{1D droplet} & $4/3 = 1.3333$ & $23/15 = 1.5333$ & $176/105 = 1.6762$\\
   \end{tabular}
    \end{ruledtabular}
    \caption{$\tau_n=\xi_{n=1}/\xi_n$ for the two OBC systems, $L^2\times M$ with $M=48$ ($L=16$) and $M=88$ ($L=24$), at several temperatures and at the critical point, read off at $r_\mathrm{min}=8$ as in Table~\ref{tab:xi}. The last row is the one-dimensional droplet-model prediction, $\sum_{k=1}^{n}1/(2k-1)$.}
    \label{tab:tau}
\end{table}
